% Generated by scripts/assemble.py; corrections belong in source/corrections.json.
\examhead{48}{东华大学 2025 年数学分析试卷}{50}
% p051-b01; source PDF page 51
\par\Needspace{5\baselineskip}
\blockanchor{p051-b01}
\majorquestion{一}计算题. 每题 8 分, 共 40 分.\par
% p051-b02; source PDF page 51
\blockanchor{p051-b02}
(1) 求极限 $\lim_{n\to \infty}n\sin (2\pi n!e)$.\par
% p051-b03; source PDF page 51
\blockanchor{p051-b03}
(2) 求极限 $\lim_{x\to 0}\left(\frac{\sin x}{x}\right)^{\frac{1}{1 - \cos x}}.$\par
% p051-b04; source PDF page 51
\blockanchor{p051-b04}
(3) 求积分  $\int_{0}^{\frac{\pi}{4}}\ln\sin(2x)\mathrm{d}x.$\par
% p051-b05; source PDF page 51
\blockanchor{p051-b05}
(4) 求曲线积分 $\oint_{L} \frac{y \, \mathrm{d}x - x \, \mathrm{d}y}{3x^{2} - 2xy + 3y^{2}}$, 其中 $L$ 为 $|x| + |y| = 1$, 沿逆时针一周.\par
% p051-b06; source PDF page 51
\blockanchor{p051-b06}
(5) 设 $\Omega = \{(x, y, z) \mid x^2 + y^2 + z^2 \leq 1, x \geq 0, y \geq 0, z \geq 0\}$, 求 $\iiint_{\Omega} \left( \frac{x^2}{a^2} + \frac{y^2}{b^2} + \frac{z^2}{c^2} \right) \mathrm{d}V$.\par
% p051-b07; source PDF page 51
\par\Needspace{5\baselineskip}
\blockanchor{p051-b07}
\majorquestion{二}(15 分) 设 $f(x)$ 在区间 $[a, b]$ 上为增函数, 且 $f(a) \geq a$, $f(b) \leq b$, 证明: 存在 $\xi \in [a, b]$, 使得 $f(\xi) = \xi$.\par
% p051-b08; source PDF page 51
\par\Needspace{5\baselineskip}
\par\Needspace{6\baselineskip}
\blockanchor{p051-b08}
\majorquestion{三}(15 分) 设函数 $f(x)$ 在区间 [0, 2024] 上连续(应该是连续可微), 且 $f(0) = 0$. 证明:\par
\nopagebreak[4]
% p051-b09; source PDF page 51
\blockanchor{p051-b09}
\displayformula{\int_ {0} ^ {2024} | f (x) f ^ {\prime} (x) | \mathrm{d} x \leq 1012 \int_ {0} ^ {2024} (f ^ {\prime} (x)) ^ {2} \mathrm{d} x.}
% p051-b10; source PDF page 51
\par\Needspace{5\baselineskip}
\par\Needspace{6\baselineskip}
\blockanchor{p051-b10}
\majorquestion{四}(15 分) 讨论函数\par
\nopagebreak[4]
% p051-b11; source PDF page 51
\blockanchor{p051-b11}
\displayformula{f (x, y) = \left\{ \begin{array}{l l} x - y + \frac {x y ^ {3}}{x ^ {2} + y ^ {4}}, & x ^ {2} + y ^ {2} \neq 0; \\ 0, & x ^ {2} + y ^ {2} = 0. \end{array} \right.}
% p051-b12; source PDF page 51
\blockanchor{p051-b12}
在原点 $(0,0)$ 是否连续，偏导数是否存在，是否可微？\par
% p051-b13; source PDF page 51
\par\Needspace{5\baselineskip}
\blockanchor{p051-b13}
\majorquestion{五}解答如下问题. 每题 15 分, 共 45 分.\par
% p051-b14; source PDF page 51
\blockanchor{p051-b14}
(1) 判断并证明级数 $\sum_{n=1}^{\infty} \sin(\sqrt{n^2 + 1}\pi)$ 的收敛性.\par
% p051-b15; source PDF page 51
\blockanchor{p051-b15}
(2) 判断并证明积分 $\int_{1}^{+\infty} \frac{y \sin(xy)}{x(1 + y^2)} \, \mathrm{d}y$ 在 (0,1) 上的一致收敛性.\par
% p051-b16; source PDF page 51
\blockanchor{p051-b16}
(3) 判断并证明函数列 $f_{n}(x) = nx(1 - x)^{n}$ 在 $[0, 1]$ 上的一致收敛性.\par
% p051-b17; source PDF page 51
\par\Needspace{5\baselineskip}
\blockanchor{p051-b17}
\majorquestion{六}(20 分) 证明: $\sum_{n=1}^{\infty} \frac{1}{n^n} = \int_0^1 x^{-x} \, \mathrm{d}x$.\par
